%% ~~~~~~~~~~~~~~~~~~~~~~~~~~~~~~~~~~~~~~~~~~~~~~~~~~~~~~~~~~~~~~~~~~ %%
%% Coastal Complexity & Resilience -- summary ~~~~~~~~~~~~~~~~~~~~~~~ %%
%% created 4/15/2022; zhr ~~~~~~~~~~~~~~~~~~~~~~~~~~~~~~~~~~~~~~~~~~~ %%
%% ~~~~~~~~~~~~~~~~~~~~~~~~~~~~~~~~~~~~~~~~~~~~~~~~~~~~~~~~~~~~~~~~~~ %%





%% preamble: load packages / document settings ~~~~~~~~~~~~~~~~~~~~~~ %%
\documentclass[11pt]{article}
\usepackage[dvipsnames]{xcolor}
\usepackage[position=top]{subcaption}
\captionsetup[subfigure]{skip=0pt}
\usepackage{float}
\floatstyle{plaintop}
\restylefloat{table}
\usepackage{graphicx}
\graphicspath{{./figs/}}
\usepackage[export]{adjustbox}
\usepackage{caption}
\usepackage{subcaption}
\usepackage{wrapfig}
\usepackage{setspace} 
\usepackage{ragged2e}
\usepackage{authblk}

\usepackage{enumitem}
\setlist[enumerate,1]{label=(\arabic*),font=\itshape}

%% biblio information 
\usepackage[
sorting=none,
backend=biber,
style=numeric,
url=false,
eprint=false,
isbn=false
]{biblatex}
\addbibresource{./biblio/references.bib}
%% end biblio 

\usepackage[margin=1in]{geometry}
\usepackage[hypertexnames=false]{hyperref}
\hypersetup{urlbordercolor=blue}
\hypersetup{linktocpage=true}
%% END preamble ~~~~~~~~~~~~~~~~~~~~~~~~~~~~~~~~~~~~~~~~~~~~~~~~~~~~~ %%

%% custom commands ~~~~~~~~~~~~~~~~~~~~~~~~~~~~~~~~~~~~~~~~~~~~~~~~~~ %%
\newcommand{\w}{.3275\textwidth} % set figure width 
\newcommand{\h}{.265\textwidth} % set figure height
\newcommand{\pt}{\vspace{-.5em}} % reduce vertical spacing 
\newcommand{\ptt}{\vspace{-.9em}} % reduce vertical spacing 
%% END custom commands ~~~~~~~~~~~~~~~~~~~~~~~~~~~~~~~~~~~~~~~~~~~~~~ %%





%% title page and title information ~~~~~~~~~~~~~~~~~~~~~~~~~~~~~~~~~ %%
\title{
research overview:
\\
Coastal Complexity \& Resilience 
\\
}
\author{
\Large{
\textbf{Seattle Aquarium} 
$\vert$ 
\textit{Conservation Programs and Partnerships}
}
\\
\large{
Zachary Randell, Ph.D., Research Scientist
\\
Shawn Larson, Ph.D., Curator of Conservation Research
\\
Erin Meyer, Ph.D., Director
}
\\
\vspace{10pt}
} 


%\textit{\textbf{
%Conservation Programs \& Partnerships, Seattle 
%Aquarium
%}}
\begin{document}
\begin{titlepage}
\maketitle
\thispagestyle{empty}
%\tableofcontents
\vfill

%% section: Overview ~~~~~~~~~~~~~~~~~~~~~~~~~~~~~~~~~~~~~~~~~~~~~~~~ %%
\section{Overview}
\label{Overview}
%% Overview ~~~~~~~~~~~~~~~~~~~~~~~~~~~~~~~~~~~~~~~~~~~~~~~~~~~~~~~~~~~~
Coastal species such as kelp, invertebrates, fishes, and marine mammals 
are increasingly threatened from climate-change related disturbances 
such as anomalous warm-water and marine-disease events. 
Unfortunately, our ability to understand the health and trends of these 
species is limited by the lack of widespread monitoring, as shallow, 
benthic ecosystems are logistically difficult to survey.   
\textit{Unmanned systems such as Remotely Operated Vehicles (ROVs) and 
Artificial Intelligence (AI) methods of analysis have the potential to 
radically alter the amount of data we can collect, analyze, and use to 
guide fishery and conservation decision-making, but thus far these 
tools have yet to be systematically incorporated into a benthic 
monitoring program.} 

In conjunction with Tribal, state, federal, and non-profit partners, Conservation Programs and Partnerships (CPP) at the Seattle Aquarium is initiating a new research project to fill this data collection and analysis need: Coastal Complexity and Resilience (CCR).
The CCR project will expand upon ongoing CPP long-term monitoring of 
sea otters along the Outer Coast and rockfish in the western Strait of 
Juan de Fuca and throughout Puget Sound, while simultaneously utilizing 
novel ROV and AI methods to establish new subtidal sentinel sites 
throughout these regions. 
In parallel---and in keeping with Seattle Aquarium values of Honoring 
Place and Learning and Improving---we have engaged with coastal Tribes, 
schools, and communities in order to share our motivation, methods, and 
findings with those livest closest to these ecosystems. 
\textit{Our ultimate objective: better understand the biological 
processes---such as the role of keystone sea otter predation---that 
provide resilience to Washington's coastal ecosystems.}
%% END section: Overview ~~~~~~~~~~~~~~~~~~~~~~~~~~~~~~~~~~~~~~~~~~~~ %%

\vspace{20pt}
\centering
\includegraphics[width=3.5cm]{logo.jpg}
\\
\end{titlepage}
%% END title page~~~~~~~~~~~~~~~~~~~~~~~~~~~~~~~~~~~~~~~~~~~~~~~~~~~~ %%





%% section: Motivation ~~~~~~~~~~~~~~~~~~~~~~~~~~~~~~~~~~~~~~~~~~~~~~ %%
\section{Motivation}
\label{Motivation}
Climate change poses a significant threat to coastal ecosystems. 
This was recently demonstrated during the 2013-2015 anomalous 
warm-water event (i.e., “The Blob”) and subsequent El Niño phase along 
the Pacific northeast: 
in northern California alone, 
more than 90\% of bull kelp was lost along 350 km of coastline, 
sea urchin density increased 60-fold, 
the sunflower star became locally extinct, 
and red abalone declined via starvation by up to 96\%, resulting in the 
closure of its productive fishery in California and Oregon (292 mt per 
year) \cite{Rogers-Bennett2019a, Gravem2021}.
Curiously, despite widespread kelp loss to the north (British Columbia) 
and south (Oregon), canopy-forming kelp along Washington's Outer Coast 
appear to have largely recovered after a relatively minor decline 
associated with the 2013-2015 anomalous climate events (Helen Berry, 
DNR; Dr. Andrew Shelton, NMFS). 
Despite experiencing widespread sea star wasting, the rapid kelp 
rebound suggests the Outer Coast may exhibit a degree of resilience, 
i.e., the capacity to recover after a disturbance.
The exact reason(s) why remain(s) unclear. 
At a regional scale (e.g., coastal Washington versus Oregon), it is 
possible the local sea otter (\textit{Enhydra lutris kenyoni}) 
population contributed to this resilience \cite{Rasher2020, Shelton2018}---yet patches of urchin barrens still developed offshore of Tatoosh and the Western Strait following the 2013-2015 disturbances. 
Environmental drivers are important for structuring kelp dynamics along 
the Outer Coast \cite{Pfister2018}---yet the same broad environmental 
conditions lead to widespread kelp loss elsewhere. 
Given the likelihood of future climate-related disturbance events, understanding the health, trends, and functioning of coastal habitats and species---all required for 
effective conservation and management---requires additional information \cite{Hughes2017}. 
%% END section: Motivation ~~~~~~~~~~~~~~~~~~~~~~~~~~~~~~~~~~~~~~~~~~ %%





%% section: Methods ~~~~~~~~~~~~~~~~~~~~~~~~~~~~~~~~~~~~~~~~~~~~~~ %%
\section{Methods}
\label{Methods}
In collaboration and partnership with
the Makah Tribe, 
National Marine Fisheries Service (NMFS), 
Olympic Coast National Marine Sanctuary (OCNMS),
Puget Sound Restoration Fund (PSRF),
Washington Dept. Fish and Wildlife (WDFW),
and with support from the 
Sea Otter Foundation and Trust,
North Pacific Coast Marine Resources Committee, and 
the Port of Seattle,
the Seattle Aquarium is launching a new project to better understand the patterns of and processing affecting benthic community structure and ecological resilience in Puget Sound and along the Outer Coast. 
%
The CRR project will utilize: 
(\textit{1}) cutting edge surveillance technology to gather imagery of 
shallow (5-40m) subtidal ecosystems including kelp forests and urchin 
barrens 
and  
(\textit{2}) machine learning algorithms to process our imagery at scale, all while 
(\textit{3}) engaging coastal communities and schools about the health of their ecosystems. 
%

\subsection{\textit{Long-term ROV monitoring of subtidal sentinel 
sites}}
Unmanned systems such as ROVs have the potential to radically expand 
the spatial extent across which we can gather information, but thus far 
their large size and expense has mostly restricted ROV usage to 
exploring deep locations unsuitable for divers \cite{Hamel2020}.
However, technological advancement has produced models that are small, 
affordable, capture high-resolution imagery, and are deployable from 
small vessels \cite{Buscher2020}. 
At present, the Seattle Aquarium has obtained, tested, and customized a 
\href{https://bluerobotics.com/}{BlueROV2} from Blue Robotics, named 
Waksa·s, a Makah word for octopus. 
We will use Waksa·s to conduct transect-based surveys along fixed 
subtidal index sites, gathering photos and video of benthic species 
such as understory and canopy-forming kelp, invertebrates such as sea 
stars, sea urchins, and abalone, and fishes such as rockfish. 
These surveys will first expand upon existing Seattle Aquarium index 
sites, and then in collaboration with our Tribal partners, we will 
establish a new network of index sites along the western Strait and 
Outer Coast. 

\subsection{\textit{AI algorithms to process ROV-imagery at scale}}
As it would not be feasible for a human to review the vast sum of 
photos and video we intend to collect, we have trained machine learning 
algorithms to extract community data from our imagery. 
Specifically, we are using \href{https://coralnet.ucsd.edu/}{CoralNet} 
\cite{Williams2019} to calculate metrics of percent-cover for 
aggregated taxa (e.g., sponges, understory algae, tunicates, as well as 
substrate type such as sand, cobble, shell debris, mud, hard 
substrate), and 
\href{https://www.star.nesdis.noaa.gov/star/documents/meetings/2020AI/presentations/202008/20200827_Dawkins.pdf}{VIAME}
 \cite{Dawkins2017} to calculate abundances of individually conspicuous 
 species (e.g., sea urchins, sea stars, abalone, fishes). 
By pairing large amounts of community data with existing habitat layers 
(e.g., from side-scan sonar derived benthic mapping), we will 
investigate the patterns, trends, and health of species across larger 
areas than SCUBA divers are logistically able to operate. 
And by regularly surveying these index sites as part of a long-term 
monitoring program, our research will directly investigate the patterns 
of and processes modifying the resilience of coastal ecosystems 
following disturbance events. 
	
\subsection{\textit{Engage with coastal communities and youth}}
Finally, in order to share our motivation, methods, and findings with 
coastal communities, we are expanded upon existing relationships 
and forging new ties along the Outer Coast. 
The Seattle Aquarium has a long relationship with the Makah, and we are 
working closely with Makah Fisheries Management (MFM) personnel in all 
aspects of this project. 
We have also established a dialogue with natural resource managers for 
the Hoh Tribe, Quileute Tribe, Quinault Indian Nation, and Samish 
Indian Nation. 
We are presently working with teachers at Neah Bay High School and the 
Quileute Tribal School to create and deliver lecture content about kelp 
forests, sea otters, climate change, and the ROV and AI tools we are 
using to evaluate coastal ecosystems.
We are also working with the Seattle Aquarium's Community Education and 
Learning department to modify one of their existing Community Science 
(CS) activities, a multi-day introduction to data collection, 
hypothesis generation, data analysis, and research presentation.
Our version of CS will utilize the ROV, AI, coding: students will 
participate in a ROV dive (e.g., in the Makah harbor for Neah Bay High 
School students),
formulate a testable hypothesis based on the imagery, 
use VIAME to annotate the images,
visualize and analyze the exported data with pre-configured scripts in 
\textit{R}-studio, and finally, present their findings to their peers. 
Experiencing ROV-based research in their own backyard, so to speak, 
along with the introduction to coding---all with a backdrop of 
evaluating ecologically pressing questions---will provide meaningful 
STEM exposure to areas of coastal Washington that may not typically 
receive much outreach attention. 
%% END section: Motivation ~~~~~~~~~~~~~~~~~~~~~~~~~~~~~~~~~~~~~~~~~~ %%





%\newpage
%section: Additional information ~~~~~~~~~~~~~~~~~~~~~~~~~~~~~~~~~~ %%
\section{Additional information}
\label{Additional information}
Click \href{https://www.youtube.com/watch?v=q3xT6H-Ufks\&t=1702s}{here} 
to see a presentation at the November 2021 Puget Sound Kelp Research 
and 
Monitoring Workgroup, given during the early stage of the CCR project. 
Furthermore, you can view lead Principal Investigator Randell being 
\href{https://www.youtube.com/watch?v=KDEc3Q2V5XQ\&ab_channel=SeaOtterFoundation\%26Trust}{interviewed}
 by Diane Tomecek, CEO and President of the Sea Otter 
Foundation and Trust (SOFT) regarding the CCR project. 
To view a recent synopsis of ROV testing and the proof-of-concept 
development of AI methods of analysis, download the report linked 
\href{https://drive.google.com/file/d/1PeGnilFcc_J5Gx2AwyMfiZZa_116hup7/view?usp=sharing}{here}.
When opened in a pdf reader, this document links to videos from the ROV.

In an effort to make research open-source and available to all, Randell 
has his Ph.D. public dissertation defense on YouTube 
\href{https://www.youtube.com/watch?v=PpVOXf6_830\&t=162s}{here}, and 
his full dissertation can be found 
\href{https://ir.library.oregonstate.edu/concern/graduate_thesis_or_dissertations/2v23w237z}{here}.
Finally, you can view more information on Randell's 
\href{https://www.researchgate.net/profile/Zachary-Randell}{ResearchGate}
 and \href{https://github.com/zhrandell}{GitHub} page. 
% END section: Additional informat 
%~~~~~~~~~~~~~~~~~~~~~~~~~~~~~~~~~~~~~~~~~~ %%





%% section: References ~~~~~~~~~~~~~~~~~~~~~~~~~~~~~~~~~~~~~~~~~~~~~~ %%
\clearpage
\section{References}
\printbibliography[heading=none]
\label{References}
%% END section: References ~~~~~~~~~~~~~~~~~~~~~~~~~~~~~~~~~~~~~~~~~~ %%





%% END document ~~~~~~~~~~~~~~~~~~~~~~~~~~~~~~~~~~~~~~~~~~~~~~~~~~~~~ %%
\end{document}
%% END document ~~~~~~~~~~~~~~~~~~~~~~~~~~~~~~~~~~~~~~~~~~~~~~~~~~~~~ %%
















